contrary to Lemma~\ref{lem:side}.
\end{proof}

\begin{lemma}[Alternating lift identity]\label{lem:alt-lift}
For every positive integer $z$ there are $m\ge1$ and
$\delta=1-(z\bmod2)$ such that
\[
  3z+1=Q_m^\delta(J(R(z))).
\]
If $R(z)>0$, one may take $m$ equal to the length of the maximal alternating
suffix of $z$; if $R(z)=0$, one takes one more than the binary length of $z$.
\end{lemma}

\begin{proof}
First suppose $r=R(z)>0$, let $k$ be the alternating-suffix length and put
$\epsilon=r\bmod2$. Then
\[
 z=2^k r+t_{k,\epsilon},\qquad
 t_{k,0}=\Big\lfloor\frac{2^k}{3}\Big\rfloor,
 \quad t_{k,1}=\Big\lfloor\frac{2^{k+1}}{3}\Big\rfloor.
\]
Since $J(r)=3r+1+\epsilon$, direct substitution gives
\[
 3t_{k,\epsilon}+1+\delta=2^k(1+\epsilon),
\]
the two parity cases being exactly the two remainders of a power of $2$
modulo $3$. Hence
\[
 3z+1+\delta=2^kJ(r),
\]
which is the required identity with $m=k$. If $R(z)=0$, the whole binary word
of $z$ is alternating and starts with $1$, so
$z=t_{k,1}$. The same calculation gives
$3z+1+\delta=2^{k+1}=2^{k+1}J(0)$.
\end{proof}

\begin{lemma}[Exceptional child-pair identity]\label{lem:exceptional-side}
Let $q\equiv1,3,12,14\pmod{16}$ and put $\ell=B(q)$.  Then for some
$m\ge1$ and $\delta\in\{0,1\}$ the two children of $A(q)$ are exactly
\[
  Q_m^\delta(J(\ell)),\qquad Q_{m+1}^\delta(J(\ell)).
\]
Consequently, if in addition $q$ is \WIN, $A(q)$ is \DRAW, and
$\ell$ is \LOSS, this adjacent pair contains a continuing \DRAW.
\end{lemma}

\begin{proof}
Write $q=16t+c$ with $c\in\{1,3,12,14\}$. The four direct exceptional
formulas are
\[
\begin{array}{c|c|c|c}
c&\ell&B(A(q))&A^2(q)\\ \hline
1&R(6t)&18t+1&36t+3\\
3&R(6t+1)&18t+4&36t+8\\
12&R(6t+4)&18t+13&36t+27\\
14&R(6t+5)&18t+16&36t+32.
\end{array}
\]
In every row put respectively $z=6t,6t+1,6t+4,6t+5$. Then
$\ell=R(z)$ and $B(A(q))=3z+1$. By Lemma~\ref{lem:alt-lift},
\[
  B(A(q))=Q_m^\delta(J(\ell))
\]
for some $m\ge1$. The last column is $2B(A(q))+\delta$, so it equals
$Q_{m+1}^\delta(J(\ell))$. These identities use no outcome hypothesis and
are precisely the two children of $A(q)$.  If $A(q)$ is \DRAW, neither
child can be \LOSS and at least one is \DRAW, proving the consequence.
\end{proof}

For a source $x$ and phase $e$, let $O(x,e)$ denote the following exact
\emph{DRAW obligation}:
\begin{enumerate}[label=(\alph*)]
\item $Q_1^e(J(x))$ is \DRAW; or
\item there is a \DRAW\ state whose two children are exactly
      $Q_2^e(J(x))$ and $Q_1^e(J(x))$.
\end{enumerate}
This definition introduces no outcome assumption beyond the displayed
\DRAW\ state.

\begin{lemma}[Hidden parent for a typed three/two frame]\label{lem:hidden-parent}
Let $a$ be positive, odd and not divisible by $3$, and suppose
\[
  a\equiv1+e\pmod3.
\]
Put $H=Q_3^e(a)$, $G=Q_2^e(a)$ and
\[
  K=\frac{2H+e-1}{3}.
\]
Then $K$ is an integer and
\[
  \Moves(K)=\{H,G\}.
\]
\end{lemma}

\begin{proof}
For $e=0$ the congruence gives $H\equiv2\pmod3$ and
$K=(2H-1)/3$; for $e=1$ it gives $H\equiv0\pmod3$ and $K=2H/3$.
These are exactly the two inverse branches of the injective map
$A(q)=\lceil3q/2\rceil$. The last three bits of $H$ are equal, so the
contracting child removes only the last bit and equals $G$.
\end{proof}

\section{Fixed-tail factor removal and entry typing}

The next two lemmas isolate the only raw entry phenomenon that is not already
an obligation or a factor frame.

\begin{lemma}[Fixed-tail factor-removal diamond]\label{lem:fixed-tail-factor}
Let $c$ be positive and odd, $e\in\{0,1\}$, and $r\ge2$. Put
\[
  X=Q_r^e(3c),\qquad G=Q_r^e(c),\qquad H=Q_{r+1}^e(c),\qquad Y=A(G).
\]
Then
\[
  \Moves(H)=\{X,Y\}.
\]
If $X$ is \DRAW, either $G$ is \DRAW, or the diamond determines a finite
routing witness together with its exact incidence to $G,H,X,Y$ and the other
child of $G$.
\end{lemma}

\begin{proof}
Since $r+1\ge3$, Lemma~\ref{lem:long-tail} gives $A(H)=X$. For $r\ge3$ it
also gives $B(H)=Q_{r-1}^e(3c)=A(G)$; for $r=2$ the boundary identity gives
the same equality. Thus $B(H)=Y$.

Because $X$ is \DRAW, $H$ is not \LOSS. If $H$ is \WIN, then its other
child $Y$ is necessarily \LOSS, and this is the required finite witness.
Assume instead that $H$ is \DRAW. Then $Y$, being a child of a DRAW, is not
LOSS. If $G$ is \LOSS, then both children of $G$, including $Y=A(G)$, are
WIN. If $G$ is \WIN, at least one child of $G$ is LOSS; since $Y$ is not LOSS,
that witness is the other child $B(G)$. These are all outcome orientations.
In every finite branch the witness and the parent-child incidence just
displayed are known exactly; no token-rank comparison is asserted at this
stage.
\end{proof}

\begin{definition}[Entry controls]\label{def:entry-controls}
A \emph{boundary entry} is an obligation $O(x,e)$ or, more generally, a
factor-free adjacent frame
\[
 \{Q_{m+1}^e(J(x)),Q_m^e(J(x))\},\qquad m\ge1,
\]
known to contain a continuing \DRAW. Its remaining tail length is recorded as
$\tau$. The typed three/two frame satisfying Lemma~\ref{lem:hidden-parent} is
a distinguished $m=2$ boundary constructor.
A \emph{loss-sibling entry} is one of the finitely many local constructors
in which a continuing \DRAW\ is accompanied by a finite \WIN/\LOSS\ routing
witness with a proved parent/child incidence: the fixed-tail diamond of Lemma~\ref{lem:fixed-tail-factor}, the exponent-one predecessor diamond of Lemma~\ref{lem:exp-one-lift}, an ordinary side diamond of Lemma~\ref{lem:side}, the exceptional side attachment of Lemma~\ref{lem:exceptional-side}, or the hidden-parent frame of Lemma~\ref{lem:hidden-parent}. The constructor records its remaining
constant-tail/factor exponent as the local counter $\tau$.
\end{definition}

\begin{lemma}[Entry typing]\label{lem:entry-typing}
Every finite branch produced by Lemma~\ref{lem:fixed-tail-factor} is a
loss-sibling entry in the sense of Definition~\ref{def:entry-controls}; no
ranked multiset is changed in order to make this classification.
\end{lemma}

\begin{proof}
The proof of Lemma~\ref{lem:fixed-tail-factor} gives the finite witness and
its exact incidence to $G,H,X,Y$ and the other child of $G$. These are
precisely the data of the first loss-sibling constructor. Classification is
therefore immediate; subsequent tail/factor normalization belongs to the
typed routing relation and is ranked there by $\tau$ and, when available,
by descendant-token comparisons.
\end{proof}

\begin{lemma}[Exponent-one lift]\label{lem:exp-one-lift}
Let $a$ be positive and odd and $k\ge1$. If
$R_k=Q_1^e(3^ka)$ is \DRAW, then
\[
  P_{k-1}=Q_2^e(3^{k-1}a)
\]
is either \DRAW\ or \WIN; in the latter case its other child is a finite
\LOSS\ witness.
\end{lemma}

\begin{proof}
The boundary recurrence gives $A(P_{k-1})=R_k$. A state with a \DRAW\ child
cannot be \LOSS. If $P_{k-1}$ is \WIN, the \DRAW\ child cannot witness the
win, so its other child is \LOSS.
\end{proof}

\begin{proposition}[Arbitrary factorful exponent-one entry]\label{prop:factorful-entry}
Let $a=J(s)$ and $k>0$. If $Q_1^e(3^ka)$ is \DRAW, then after finitely many
outcome-compatible local steps one obtains either a typed entry of
Lemma~\ref{lem:entry-typing}, or the factor-free exponent-two state
\[
  Q_2^e(a)\text{ \DRAW}.
\]
\end{proposition}

\begin{proof}
Apply Lemma~\ref{lem:exp-one-lift}. Its finite branch is a loss-sibling entry by Definition~\ref{def:entry-controls}. Otherwise $Q_2^e(3^{k-1}a)$ is \DRAW. Apply
Lemma~\ref{lem:fixed-tail-factor} with $r=2$ repeatedly. Whenever its finite
branch occurs, Lemma~\ref{lem:entry-typing} supplies a typed entry. If it never
occurs, one factor of $3$ is removed at the same tail exponent at each step,
so after $k-1$ steps the displayed factor-free exponent-two state is reached.
\end{proof}

\begin{proposition}[Complete exponent-one/two boundary reduction]\label{prop:boundary-reduction}
Let $s\ge0$, $k\ge0$, $e\in\{0,1\}$, and let
\[
  Q_r^e(3^kJ(s)),\qquad r\in\{1,2\},
\]
be \DRAW. For $s>0$, after finitely many local steps at least one of the following occurs: a typed entry is obtained, the retained source strictly
decreases, or the factor-free exponent-two state $Q_2^e(J(s))$ is \DRAW.
For $s=0$ the corresponding statement is Lemma~\ref{lem:zero-source}.
\end{proposition}

\begin{proof}
Assume $s>0$. There are four cases. If $r=1,k=0$, the state is the first
alternative of the obligation $O(s,e)$ and is already a boundary entry. If
$r=1,k>0$, use Proposition~\ref{prop:factorful-entry}. If $r=2,k=0$, we are
at the displayed factor-free exponent-two state. Finally, if $r=2,k>0$, apply
Lemma~\ref{lem:fixed-tail-factor} repeatedly with $r=2$: a finite branch is a
loss-sibling entry by Definition~\ref{def:entry-controls}, while otherwise
one factor of $3$ is removed at the same exponent at each step and after $k$
steps the factor-free exponent-two state is reached. These cases are disjoint
and exhaustive.
\end{proof}

\section{Factor-free exponent-two base entry}

Let $x>0$, $a=J(x)$ and
\[
  P=Q_1^e(a),\quad U=Q_2^e(a),\quad
  b=B(P)=B(U),\quad F=A(U)=Q_1^e(3a),\quad g=1-e.
\]
Assume $U$ is \DRAW\ and $x$ is the least coefficient source of any \DRAW.  This is a \emph{factor-free} coefficient $J(x)$.  The common child $b$ cannot be $0$, because $0$ is \LOSS\ and a \DRAW\ has no \LOSS\ child.  By Lemma~\ref{lem:boundary-coeff}, the canonical coefficient of $b$ is strictly below $J(x)$.  After deleting any factor of $3$ its three-free coefficient is smaller still, hence it equals $J(t)$ with $t<x$ because $J$ is strictly increasing.  Therefore the state source of $b$ is below $x$, and $b$ is not \DRAW.  No such inference is made for a coefficient $3^kJ(x)$ with $k>0$.
As a child of a \DRAW\ it is not \LOSS, so
\[
  b\text{ is \WIN},\qquad F\text{ is \DRAW}.
\]
\begin{lemma}[First-factor anchor identity]\label{lem:first-factor}
For the states above,
\begin{equation}\label{eq:first-factor}
  9J(x)+1-2e=2^vJ(b),\qquad v\ge1.
\end{equation}
\end{lemma}

\begin{proof}
Put $a=J(x)$. If $e=0$, then $P=2a$, $A(P)=3a$ and $b=R(3a)$.
Apply the first identity of Lemma~\ref{lem:coeff-pullback} to the odd integer
$w=3a$:
\[
 \odd(9a+1)=\kappa(9a)=J(R(3a))=J(b).
\]
If $e=1$, then $P=2a-1$, $A(P)=3a-1$ and $b=R(3a-1)$. The second identity
of Lemma~\ref{lem:coeff-pullback}, again with $w=3a$, gives
\[
 \odd(9a-1)=J(R(6a-1)).
\]
By the boundary identity of Lemma~\ref{lem:long-tail},
$R(6a-1)=R(3a-1)=b$. Hence in both phases the odd part of
$9a+1-2e$ is exactly $J(b)$, which proves \eqref{eq:first-factor}.
\end{proof}

\begin{proposition}[Base-entry attachment]\label{prop:base-entry}
Every factor-free exponent-two \DRAW\ at a globally minimal source has a
finite continuation which either strictly lowers the retained source or
produces a boundary/loss-sibling typed entry.
\end{proposition}

\begin{proof}
Suppose first that $e=1-\alpha(x)$. The first source-boundary numerator is
divisible by $4$, while
\[
9J(x)+1-2e=3(3J(x)+1-2e)-2(1-2e)\equiv2\pmod4.
\]
Thus $v=1$. Writing $u=A(x)$ gives $b=3u+1$. The two children of the \DRAW\
state $F$ are
\[
  T=Q_1^g(J(b)),\qquad C=R(T),
\]
and direct appended-bit deletion gives
$C\in\{A(b),B(b)\}$. If $T$ is \DRAW, this is the boundary entry
$O(b,g)$. If $C$ is \DRAW, the other child $\ell$ of the \WIN\ state $b$ is a
finite \LOSS\ witness. If $b$ is ordinary, this is the ordinary-side
loss-sibling constructor. If $b$ is exceptional and $C=A(b)$,
Lemma~\ref{lem:exceptional-side} gives an adjacent \DRAW\ frame over
$\ell$, hence a boundary/loss-sibling entry. It remains only the exceptional
orientation $C=B(b)$. In that orientation
$b\le(9x+5)/2$ and the alternating suffix of $A(b)$ has length at least three,
so
\[
 C=B(b)\le\frac{A(b)}8\le\frac{3b+1}{16}.
\]
Lemma~\ref{lem:source-bound} then yields
\[
 \rho(C)\le\frac{C-1}{6}<\frac{3b+1}{96}
 \le\frac{27x+17}{192}<x,
\]
a strict retained-source exit.

Now suppose $e=\alpha(x)$. Lemma~\ref{lem:source-boundary} gives
$3J(x)+1-2e=2J(A(x))$, hence $v\ge2$. If $v\ge4$, then
\[
16J(b)\le9J(x)+1\le27x+19,
\]
and $J(b)\ge3b+1$ implies
\[
  b\le\frac{9x+1}{16}<x.
\]
Thus the continuing \DRAW\ has smaller retained source. If $v=2$, the two
children of $F$ are exactly $Q_2^g(J(b))$ and $Q_1^g(J(b))$, which is the
second form of $O(b,g)$. If $v=3$, they are the typed three/two frame
$Q_3^g(J(b)),Q_2^g(J(b))$. Reducing
\eqref{eq:first-factor} modulo $3$ gives
\[
  J(b)\equiv1+g\pmod3,
\]
so Lemma~\ref{lem:hidden-parent} applies. The hidden parent is a boundary
entry; its finite side, if one is exposed, is a loss-sibling entry. Hence no
untyped base row remains.
\end{proof}

\section{The source-zero branch}

For $k\ge0$, $r\ge1$ and $e\in\{0,1\}$ put
\[
  S_{k,r}^e=Q_r^e(3^k).
\]
These are precisely the constant-tail states with coefficient source zero.

\begin{lemma}[Zero-source boundary normalization]\label{lem:zero-source}
If a source-zero \DRAW\ exists, then after finitely many outcome-compatible
moves one reaches a typed boundary or loss-sibling entry.  No ranked token is
installed and no control reset is used before this finite normalization is
complete.
\end{lemma}

\begin{proof}
For $r\ge3$, Lemma~\ref{lem:long-tail} gives
\[
  \Moves(S_{k,r}^e)=\{S_{k+1,r-1}^e,S_{k+1,r-2}^e\}.
\]
Following a \DRAW\ child strictly lowers $r$, so a \DRAW\ with $r\in\{1,2\}$
is reached after finitely many moves. If $k=0$ already at that boundary, the
four possibilities are explicit:
\[
 Q_1^1(1)=1\ \WIN,\quad Q_2^1(1)=3\ \WIN,\quad
 Q_1^0(1)=2\ \LOSS,\quad Q_2^0(1)=4\ \LOSS.
\]
Hence a boundary \DRAW\ must have accumulated at least one factor of $3$.

If the reached boundary has tail exponent one, put $n=k+1$.  If it has
tail exponent two, its children are the common raw child and
$S_{k+1,1}^e$.  If the raw child is \DRAW, again put $n=k+1$; otherwise the
raw child is \WIN\ and forces $S_{k+1,1}^e$ to be \DRAW, in which case put
$n=k+2$.  Thus in all cases one obtains a boundary \DRAW\ configuration whose
raw/signed exits are governed, for some $n\ge1$, by
\[
  3^n+1-2e=2^jJ(y).
\]
The exact $2$-adic valuation is
\[
\begin{array}{c|cc}
 &n\text{ odd}&n\text{ even}\\ \hline
 e=0&j=2&j=1\\
 e=1&j=1&j=2+v_2(n).
\end{array}
\]
For $j\ge2$ the raw and signed exits form the adjacent factor-free frame
\[
  Q_{j-1}^{1-e}(J(y)),\qquad Q_j^{1-e}(J(y)),
\]
containing a continuing \DRAW; by Definition~\ref{def:entry-controls} this is a boundary entry, with its finite tail length recorded in $\tau$. For $j=1$ one has
\[
  y=\frac{3^{n-1}-1}{2}.
\]
When $e=0$ the raw exit is $A(y)$; when $e=1$ it is $B(y)$. The raw and
signed states are the two exits of the same boundary \DRAW\ configuration.
If the signed exit is \DRAW, it is already the exponent-one form of an
obligation. If the raw exit is the continuing \DRAW\ and the signed exit is
finite, that finite state is necessarily \WIN; choosing a canonical \LOSS\
child of it supplies an exact finite routing witness, hence a loss-sibling
entry. If both exits are \DRAW, use the signed obligation. Thus every row
reaches a typed entry after a finite pre-entry computation; no impossible
``source below zero'' comparison is used.
\end{proof}

\begin{lemma}[$A$-selecting side diamond]\label{lem:a-side}
Let $x>0$, $e=\alpha(x)$, $w=Q_1^e(J(x))$, and $y=A(x)$. Then
\[
  B(w)\in\{A(y),B(y)\}.
\]
\end{lemma}

\begin{proof}
The source-boundary identity gives
$A(w)=Q_1^{1-e}(J(y))=4A(y)+1+e$. For $e=0$ the appended bits are $01$;
for $e=1$ they are $10$. If the first appended bit repeats the last bit of
$A(y)$, deleting the alternating suffix leaves $A(y)$; otherwise the deletion
continues through the maximal alternating suffix of $A(y)$ and leaves
$B(y)$.
\end{proof}

\begin{lemma}[Universal $B$-transfer diamond]\label{lem:b-diamond}
Let $x>0$, let $e=1-\alpha(x)$, put $a=J(x)$ and $g=1-e$, and factor
